\chapter{Notation and Conventions}\label{ch:notation}

The tomes were written to be read on their own, and each one states its conventions where it first needs them. Most of
them are shared; this chapter collects them, and points out the few places where the tomes differ.

\section{Numbers and sets}

\begin{center}\small
\begin{tabular}{@{}lp{0.62\linewidth}@{}}
$\N$, $\Z$, $\Q$, $\R$, $\C$ & the natural numbers $\{0,1,2,\dots\}$, the integers, the rationals, the reals, the complex numbers\\
$\mathbb F_q$ & the finite field with $q$ elements\\
$\Z/n\Z$, $(\Z/n\Z)^*$ & the integers modulo $n$, and its group of units\\
$p$, $q$ & in number theory, $p$ always denotes a prime\\
$\lfloor x\rfloor$, $\{x\}$, $\|x\|$ & the integer part, the fractional part, the distance to the nearest integer (for vectors, $\|x\|$ is a norm)\\
$\mathbf 1_A$ & the indicator function of a set or an event $A$\\
$\gcd$, $\operatorname{lcm}$ & greatest common divisor, least common multiple\\
$\operatorname{Re}$, $\operatorname{Im}$ & real and imaginary parts
\end{tabular}
\end{center}

Whether $0\in\N$ is the one convention of this table that the literature does not agree on; where it matters, the
tome says so.

\section{Arithmetic functions}

\begin{center}\small
\begin{tabular}{@{}lp{0.62\linewidth}@{}}
$\pi(x)$ & the number of primes $p\le x$\\
$\theta(x)$, $\psi(x)$ & Chebyshev's functions $\sum_{p\le x}\log p$ and $\sum_{n\le x}\Lambda(n)$\\
$\Lambda(n)$ & von Mangoldt's function: $\log p$ if $n$ is a power of the prime $p$, and $0$ otherwise\\
$\mu(n)$, $\lambda(n)$ & the M\"obius and Liouville functions\\
$\varphi(n)$ & Euler's totient\\
$\sigma(n)$, $d(n)$ or $\tau(n)$ & the sum and the number of the divisors of $n$\\
$\omega(n)$, $\Omega(n)$ & the number of prime factors of $n$, without and with multiplicity\\
$\zeta(s)$, $L(s,\chi)$ & the Riemann zeta function and the Dirichlet $L$-function of a character $\chi$\\
$\operatorname{li}(x)$, $\operatorname{Li}(x)$ & the logarithmic integrals $\int_0^x\frac{\dd t}{\log t}$ (principal value) and $\int_2^x\frac{\dd t}{\log t}$\\
$\operatorname{rad}(n)$, $\operatorname{ord}_p(a)$ & the product of the distinct primes dividing $n$, and the order of $a$ modulo $p$
\end{tabular}
\end{center}

\section{Analysis and asymptotics}

\begin{center}\small
\begin{tabular}{@{}lp{0.62\linewidth}@{}}
$\log$ & the natural logarithm in number theory and analysis; information theory and the computing track also use
$\log_2$, and say so\\
$f=O(g)$, $f\ll g$ & $|f|\le Cg$ for a constant $C$ (in a stated range)\\
$f=o(g)$ & $f/g\to0$\\
$f\sim g$ & $f/g\to1$\\
$f\asymp g$ & $f\ll g$ and $g\ll f$\\
$\dd x$ & the differential in an integral, written upright\\
$\Gamma(s)$ & Euler's Gamma function
\end{tabular}
\end{center}

\section{Algebra}

\begin{center}\small
\begin{tabular}{@{}lp{0.62\linewidth}@{}}
$S_n$ or $\mathfrak S_n$ & the symmetric group on $n$ letters; both spellings occur, with the same meaning\\
$A_n$ & the alternating group\\
$\operatorname{Gal}(L/K)$ & the Galois group of a Galois extension $L/K$\\
$\mathcal O_K$ & the ring of integers of a number field $K$\\
$\operatorname{Hom}$, $\ker$, $\operatorname{im}$ & homomorphisms, kernel, image\\
$\det$, $\operatorname{tr}$, $\operatorname{rank}$ & determinant, trace and rank of a matrix or a linear map\\
$\mathrm{GL}_n$, $\mathrm{SL}_n$ & the general and special linear groups
\end{tabular}
\end{center}

\section{Probability}

\begin{center}\small
\begin{tabular}{@{}lp{0.62\linewidth}@{}}
$\mathbb P$, $\mathbb E$ & probability and expectation\\
$\operatorname{Var}$, $\operatorname{Cov}$ & variance and covariance\\
$X_n\to X$ & the mode of convergence (almost sure, in probability, in law) is written over the arrow or said in words
\end{tabular}
\end{center}

\section{How the books are organised}

Every tome follows the same layout, so that a reader who knows one knows them all.
\begin{itemize}
\item Theorems, propositions, lemmas, corollaries, definitions and examples share one counter per chapter: Theorem~4.3
is the third numbered statement of Chapter~4, whatever its kind. Exercises have their own counter, also per chapter.
Remarks are not numbered; they quote deeper results, with a reference, without proof.
\item A reference to another tome always names it: ``Tome~2, Theorem~10.6''. Without a tome, the reference is to the
same book.
\item Exercises are graded: $\star$ warm-up, $\star\star$ for a good afternoon, $\star\star\star$ gems. Each has a hint,
and a solution that starts with \emph{How to find it}, the idea, before the \emph{Write-up}.
\item ``Pause and think'' boxes ask a question before a key result; the answer follows in the text.
\item Every chapter ends with sources and further reading, and comes with a Python notebook, exercises first and
solutions after.
\end{itemize}

\begin{think}
Look up Dirichlet in the index of Chapter~\ref{ch:index}. How many results carry his name, and in how many tomes?
Why would the same theorem be proved in two different tomes?
\end{think}
